\section{Core Setups I: Bounce, Breakout, and Failed Continuation}

\subsection{Bounce from visible liquidity or clustered passive interest}

\subsubsection{What the setup is}

The bounce trade is the clearest and most explicitly favored setup in the transcript set. The
basic form is simple: price approaches visible passive liquidity or a localized cluster of
passive participants, fails to trade through it cleanly, and reverses enough to deliver a
favorable short-term reward relative to a very local stop.

The materials treat bounce setups as the most educational and often the most attractive early
career trades because they are naturally anchored to something concrete. The trader is not
guessing in empty space. The trader is reacting to a specific local obstacle: one density,
multiple aligned densities, an iceberg, or a visible participant whose presence tightly defines
where the idea is wrong.

\subsubsection{Why it is believed to work}

The edge hypothesis is not that every large order creates a reversal. The edge hypothesis is
that a sufficiently meaningful passive participant can temporarily overpower aggressive flow,
especially when that participant appears at a location that already makes sense on the chart or
after a one-sided local move into the level. If the passive side is stronger than immediate
market pressure, price must either pause or reverse.

The transcript logic also adds a second layer: even when the passive participant is not large
in absolute terms, it can still matter if the market is stretched or the book is unusually
thin. The spring setup is an extreme example of this principle, but the same logic applies in
ordinary bounce trades as well.

\subsubsection{What to watch}

\paragraph{In the order book.}
Look for large visible size, multiple nearby densities, refreshing size, or an iceberg-like
refresh pattern. The most attractive cases are those where the stop can be placed just beyond
the defending participant rather than in open space.

\paragraph{In the tape.}
Watch whether aggressive prints hit the level and fail to move price. That suggests absorption.
If aggressive prints arrive and the level immediately thins, disappears, or is crossed without
friction, the bounce thesis weakens sharply.

\paragraph{In price behavior.}
A clean bounce usually has one of two shapes: either price approaches the level and stalls
before crossing, or it briefly pokes through, finds no follow-through, and snaps back. The
materials are clear that a short fake penetration does not automatically invalidate the idea.
What matters is whether there is immediate continuation through the defending participant.

\subsubsection{When not to enter}

Do not treat density alone as sufficient. Avoid bounce entries when the density is far from any
meaningful structural location, when the tape is overwhelmingly one-directional against the
level, when the displayed size is being pulled around constantly, or when the instrument is so
empty that local depth is not trustworthy. Also avoid confusing \emph{a large number on screen}
with \emph{a participant who is actually defending that level}.

\subsubsection{Stop, exit, and failure logic}

The main attraction of the bounce is the short stop. The stop is typically placed just beyond
the defending participant, the second supporting density, or the point where the level should
have held but clearly did not. Partial exits are typically justified into the first reaction
away from the level, while runners can be left for the next nearby structural area. The
failure mode is straightforward: if the limiting participant is gone or decisively consumed,
remaining in the trade becomes narrative attachment rather than disciplined execution.

\subsubsection{Clean versus unclean examples}

Clean examples have a visible reason, a readable stop, supporting context, and either
absorption or at least a lack of continuation through the level. Unclean examples are those
where the trader wants a bounce because the move ``feels extended'' but cannot point to the
actual participant or structure that should stop price.

\subsubsection{Algorithmic translation}

Possible observable features include:

\begin{itemize}[leftmargin=1.4em]
  \item displayed size at the target price as a multiple of rolling median queue size,
  \item distance from best bid / ask to the defending price,
  \item net aggressive volume into the level over the last N seconds,
  \item short-term price displacement after first contact,
  \item whether displayed size stays, refreshes, or disappears after first impact,
  \item contextual flags such as gap open, strong trend day, or proximity to prior chart level.
\end{itemize}

What remains fuzzy is the trader's subjective evaluation of whether the participant ``looks
real'' and whether the overall market feels easy or heavy. Those judgments may later require
either richer hand-engineered context or learned filters.

\subsection{Breakout through density or structural level}

\subsubsection{What the setup is}

The breakout setup is described as materially harder than the bounce. Price leans into a level
or density, pressure builds, the level gives way, and the trader attempts to join continuation
through the break. In the transcript worldview this is a high-skill trade because the trader is
often entering into a partially empty book with less obvious support behind the position.

\subsubsection{Why it is believed to work}

A valid breakout is expected to work when passive liquidity is exhausted or pulled and the
level's failure triggers additional buying or selling pressure from stops, late participants,
and momentum traders. The edge hypothesis is therefore fuel-driven continuation, not merely a
geometric line cross.

\subsubsection{What to watch}

\paragraph{In the order book.}
Compression toward the level, repeated tests, thinning liquidity on the far side, or evidence
that the defending size is smaller or weaker than it first appeared.

\paragraph{In the tape.}
Urgent prints into the level, a burst of marketable flow as price crosses, and immediate
follow-through instead of hesitation. The tape matters more here than in a plain bounce because
breakout confirmation is about acceleration.

\paragraph{In price behavior.}
The break should not just occur; it should continue. A good breakout often shows fast
acceptance above or below the level. If price crosses and then instantly hesitates in a dead
book, the trader must become suspicious immediately.

\subsubsection{Why the setup is dangerous}

The transcripts are unusually direct on this point: most breakout trades do not have a clean
stop sitting right behind them. If the trader is wrong, the position can be ``smeared'' by an
empty book. This is why breakout trading is framed as a later-skill setup requiring immediate
loss acceptance, fast reaction, and emotional discipline. A trader who cannot scratch quickly
should not be leaning on this setup heavily.

\subsubsection{Invalidation and exit logic}

The invalidation is behavioral, not just geometric. If the level breaks but there are no stops,
no follow-through, and no continuing urgency, the breakout thesis degrades almost instantly.
The trader is expected to exit fast rather than hope. Partial profit-taking makes sense into the
first acceleration burst, but only if the move is already working. Breakout management is not
about giving endless room to be right eventually; it is about being paid quickly when the read
is correct.

\subsubsection{Algorithmic translation}

Possible components include:

\begin{itemize}[leftmargin=1.4em]
  \item number of recent tests of the same level,
  \item compression metrics such as shallower pullbacks or decreasing time away from the level,
  \item aggressive flow burst at and just after the break,
  \item change in displayed far-side liquidity after the level gives way,
  \item first-second post-break excursion, retrace size, and acceptance duration.
\end{itemize}

The hard part is separating genuine continuation from noise around obvious levels. The tapes
make that possible in principle, but only if live capture preserves timing and sequencing
accurately.

\subsection{False breakout and failed continuation}

\subsubsection{What the setup is}

The false breakout is the mirror image of the breakout. A level appears to give way, the crowd
commits, but continuation fails and price returns back through the trigger area. In the logic
of the materials, this is one of the most informative events because it reveals that apparent
strength was not real strength.

\subsubsection{Why it is believed to work}

If a level attracts breakout participation and still cannot hold, someone stronger is likely on
the other side. Failed continuation also means the move has just trapped the most emotionally
committed participants. Their exits can become the fuel for the reversal.

\subsubsection{How to observe it}

Watch for a breach without sustained follow-through, quick return into the prior range,
disappearing urgency on the tape, or evidence that the move into the break was mostly a stop
run rather than genuine acceptance. In practice, this setup often shares the same ingredients as
breakout analysis, but interpreted after the move proves itself false.

\subsubsection{Algorithmic translation}

Measure:

\begin{itemize}[leftmargin=1.4em]
  \item time spent beyond the broken level,
  \item maximum excursion beyond the level before re-entry,
  \item aggressive flow during the break versus during the reversal,
  \item recovery back through the trigger zone,
  \item whether the previous liquidity that was supposed to disappear actually reappears.
\end{itemize}

This setup is an especially good candidate for later ML filtering because ``healthy breakout''
versus ``failed breakout'' often depends on short-lived interactions among depth, trade flow,
and context rather than a single threshold.
